% engines: lualatex, pdflatex
% expect-warnings: 0
% expect-errors: 0
\documentclass{article}
\usepackage[margin=1cm]{geometry}
\usepackage{../../latex/pgf-go}
\newcommand\case[1]{\par\medskip\noindent\textbf{#1}\par\nopagebreak}
\pagestyle{empty}
\begin{document}

\case{All shapes on empty points}
\begin{goban}[7]
	\marks[circle]{A1}\marks[triangle]{B1}\marks[square]{C1}\marks[cross]{D1}
	\marks[filled square]{E1}\marks[filled circle]{F1}\marks[hatched lines]{G1}
	\marks[t-circle]{A3}\marks[t-square]{B3}
	\marks[custom = {!}]{E3}\marks[custom = {AB}]{F3}
\end{goban}

\case{Marks without a player are neutral; on stones use b/w (B2 black, D2 white)}
\begin{goban}[5]
	\stones[b]{B2}\stones[w]{D2}
	\marks[b, triangle]{B2}\marks[w, triangle]{D2}\marks[triangle]{C4}
\end{goban}

\case{The n shorthand is neutral (grey filled circle)}
\begin{goban}[5]\marks[n, filled circle]{C3}\end{goban}

\case{Sequences: 1-3, from 5, Roman, and the sequence/format alias}
\begin{goban}[5]
	\marks[sequence]{A1,B1,C1}
	\marks[sequence = {from = 5}]{A3,B3}
	\marks[label = {format = Roman}, sequence]{A5,B5}
	\marks[sequence = {format = alph}]{D5,E5}
\end{goban}

\case{Mark options: red circle line, coloured label}
\begin{goban}[5]
	\marks[circle, line = {color = red}]{B2}
	\marks[label = {color = blue}, sequence]{D2,D3}
\end{goban}

\case{Fragile content in custom marks}
\begin{goban}[5]\marks[custom = \textbf{A}]{B2}\stone[w, mark = {custom = \emph{x}}]{D2}\end{goban}

\case{\textbackslash mark is an alias of \textbackslash marks inside goban}
\begin{goban}[5]\mark[circle]{C3}\end{goban}

\end{document}
